% Principios de dos acelerómetros: piezoeléctrico de cortadura con electrónica integrada (IEPE) y MEMS capacitivo diferencial.
\documentclass[tikz,border=4pt]{standalone}
\usepackage{resonancia}
\begin{document}
\begin{tikzpicture}[font=\sffamily\small, nota/.style={font=\sffamily\footnotesize, inner sep=1.5pt}]
  % =================== (a) piezoeléctrico ===================
  \node[font=\sffamily\bfseries\small, anchor=west] at (-0.8,5.1) {(a) Piezoeléctrico con electrónica integrada};
  \fill[apoyo] (0,0) rectangle (4,0.35);
  \draw[trazo] (0,0) rectangle (4,0.35);
  \node[nota, anchor=west] at (4.15,0.18) {base fijada a la estructura};
  \draw[trazo, fill=rVerde!14] (0.9,0.35) rectangle (3.1,0.95);
  \node[nota] at (2,0.65) {cerámica o cristal};
  \draw[masa, minimum width=2.6cm] (0.7,0.95) rectangle (3.3,2.05);
  \node at (2,1.5) {masa sísmica $m$};
  \draw[fuerza] (-0.6,0.2) -- (-0.6,1.3) node[midway, left, nota, color=rNaranja] {$a(t)$};
  \draw[<->, trazo, color=rVerde] (3.55,0.95) -- (3.55,1.95);
  \node[nota, anchor=west, color=rVerde, align=left] at (3.65,1.45) {$F=m\,a$ sobre\\la cerámica};
  % electrónica integrada
  \draw[trazo, fill=rAzul!8, rounded corners=2pt] (0.3,2.45) rectangle (3.7,3.45);
  \node[nota, align=center] at (2,2.95) {amplificador integrado\\alta $\to$ baja impedancia};
  \draw[trazo, color=rVerde] (0.9,0.65) -- (0.5,0.65) -- (0.5,2.45);
  \draw[flecha] (3.7,2.95) -- (4.6,2.95) node[right, nota, align=left] {salida: tensión\\con nivel de\\continua (bias)};
  \node[nota, anchor=west, align=left] at (-0.8,4.3) {Dos hilos: corriente constante hacia el sensor\\y señal de vuelta por el mismo cable};
  % =================== (b) MEMS capacitivo ===================
  \begin{scope}[xshift=9.6cm]
  \node[font=\sffamily\bfseries\small, anchor=west] at (-0.4,5.1) {(b) MEMS capacitivo diferencial};
  \fill[apoyo] (-0.35,0.2) rectangle (0,1.0); \draw[trazo] (0,0.2) -- (0,1.0);
  \fill[apoyo] (4.0,0.2) rectangle (4.35,1.0); \draw[trazo] (4.0,0.2) -- (4.0,1.0);
  \draw[masa, minimum width=2.4cm] (0.8,0.2) rectangle (3.2,0.9);
  \node at (2,0.55) {masa móvil};
  \draw[decorate, decoration={zigzag, segment length=2.4mm, amplitude=1.4mm, pre length=2mm, post length=2mm}, trazo] (0,0.55) -- (0.8,0.55);
  \draw[decorate, decoration={zigzag, segment length=2.4mm, amplitude=1.4mm, pre length=2mm, post length=2mm}, trazo] (3.2,0.55) -- (4.0,0.55);
  % dedos móviles (negro, unidos a la masa)
  \foreach \x in {1.15,2.0,2.85} { \draw[line width=1.8pt, color=rTinta] (\x,0.9) -- (\x,2.05); }
  \draw[trazo, fill=rGrisClaro] (0.6,2.25) rectangle (3.4,2.45);
  % dedos fijos: C1 azul continuo, C2 naranja discontinuo
  \foreach \x in {0.95,1.8,2.65} { \draw[line width=1.5pt, color=rAzul] (\x,2.25) -- (\x,1.2); }
  \foreach \x in {1.35,2.2,3.05} { \draw[line width=1.5pt, color=rNaranja, dash pattern=on 4pt off 2pt] (\x,2.25) -- (\x,1.2); }
  \node[nota, color=rAzul, anchor=east] at (0.5,1.75) {$C_1$};
  \node[nota, color=rNaranja, anchor=west] at (3.5,1.75) {$C_2$};
  \node[nota, anchor=west, align=left] at (-0.4,3.0) {negro: dedos móviles; azul y naranja:\\dedos fijos (dos condensadores)};
  \draw[<->, trazo, color=rGris] (1.3,-0.35) -- (2.7,-0.35) node[midway, below=3pt, nota, color=rGris] {desplazamiento $x$};
  \node[nota, anchor=west, align=left] at (-0.4,4.3) {con $a$, la masa se desplaza $x$: un condensador\\sube y el otro baja; se mide su diferencia};
  \end{scope}
  \node[nota, color=rGris, anchor=north west, align=left] at (-0.8,-0.75) {Esquemas cualitativos, no de un modelo ni de un fabricante concretos.};
\end{tikzpicture}
\end{document}
